% =====================================================================
%  Harald Høffding, LEADING IDEAS IN THE NINETEENTH CENTURY.  Copenhagen: Gyldendalske Boghandel, Nordisk Forlag, 1920.
%  (Ledende Tanker i det nittende Aarhundrede.)
%  ENGLISH TRANSLATION of an EXTRACT, printed pp. 28-32, from transcription.tex (Danish), made in the main session
%  2026-10-23.  Conventions: TRANSLATION-PLAYBOOK.md.  Page numbers of the 1920 print in the margin.
%
%  TERMS.  »Sætningen om Energiens Bestaaen« -> "the principle of the conservation of energy" (Bestaaen, literally
%  "continuance", is Høffding's word for conservation throughout: »Bevægelsens Bestaaen« -> "the conservation of motion",
%  »Stoffets og Kraftens Bestaaen« -> "the conservation of matter and force").  »Sætning« -> "principle" (or "proposition"
%  where it is a statement, not a law); »Grundsætning« -> "maxim".  »Kraft« -> "force", kept in Leibniz's sense where
%  Høffding glosses it as energy.  »Erfaringsbekræftelse« -> "confirmation by experience".  »Omsætning« -> "conversion".
%  »Horizont« -> "horizon".  The title: "Leading Ideas" for »Ledende Tanker« (literally "leading thoughts") -- not yet
%  approved by Hans.
%
%  QUOTATIONS.  Christiansen (p. 30-31) is quoted in Danish and is translated from Høffding's Danish.  Woodward (p. 31) is
%  quoted in Danish translation from an English address (R. S. Woodward, "The Unity of Physical Science", Congress of Arts
%  and Science, St. Louis 1904); the English original was not checked, so the passage is translated back from Høffding's
%  Danish and marked at the site.  The print's misprints (»Zeemann«, »Céliste«) are kept as printed, as in the Danish.
% =====================================================================
\documentclass[12pt,a4paper,openany]{book}
\usepackage[utf8]{inputenc}
\usepackage[T1]{fontenc}
\usepackage{libertinus}
\usepackage{libertinust1math}
\usepackage[english]{babel}
\usepackage[protrusion=true,expansion=true]{microtype}
\usepackage{setspace}
\usepackage[margin=1.3in]{geometry}
\usepackage{fancyhdr}
\usepackage{hyperref}
\hypersetup{pdftitle={Leading Ideas in the Nineteenth Century (pp. 28--32)}, pdfauthor={Harald Høffding}, hidelinks}
\setlength{\headheight}{15pt}
\pagestyle{fancy}
\fancyhf{}
\fancyhead[LE,RO]{\thepage}
\fancyhead[C]{\textit{Høffding: Leading Ideas in the Nineteenth Century (1920)}}
\setstretch{1.3}
\usepackage{marginnote}
\newcommand{\opage}[1]{\marginnote{\footnotesize\textit{[#1]}}}
\newcommand{\sectionhead}[2]{\par\bigskip\phantomsection\begin{center}#1 \textit{#2}\end{center}\medskip\ignorespaces}

\begin{document}

\begin{center}
{\Large Harald Høffding}\\[0.6em]
{\LARGE Leading Ideas in the Nineteenth Century}\\[0.3em]
{\large (\textit{Ledende Tanker i det nittende Aarhundrede})}\\[0.6em]
Copenhagen: Gyldendalske Boghandel, Nordisk Forlag, 1920\\[1.2em]
{\small Extract: pp.~28--32. The text begins and ends mid-sentence.}\\[0.6em]
{\small Translated from the Danish}
\end{center}
\bigskip

\noindent[\dots]
% --- p. 28 ---
% running head: LEIBNIZ OG ROBERT MAYER (Leibniz and Robert Mayer)
\opage{28}into one with the relation between ground and consequent, will not be discussed here.
But for a historical consideration it is of great interest to see that
the discovery of the conservation of energy was a continuation and confirmation by experience
of a principle that had already been put forward in the seventeenth century.
Descartes put forward a law of the conservation of motion,
according to which no motion came to nothing and no motion arose
from nothing; there was always the same sum of motion
in the universe, only it was distributed in different ways. Leibniz criticized
this doctrine, maintaining that according to it every transition
from motion to rest and from rest to motion within the universe
would be incomprehensible. There is, he further maintains, no absolute
opposition between rest and motion; rest is arrested motion,
due to an equilibrium between forces (tendencies), and motion
is released anew as soon as the equilibrium is removed. It is thus not motion
but its possibility, ``force'', that is conserved, and while there
is not always the same sum of motion in the universe, there is the
same quantity of ``force'' (energy, in the modern expression). Leibniz
laid such great weight on this principle that he even called it
the foundation of all the laws of nature (\textit{fundamentum naturæ legum}). In
the science of nature of the following period, this principle receded for various reasons,
although it has again and again been at work as
a tacit presupposition. When it now came forward again, it was not as
a proposition of pure mechanics, but as a principle asserting
a definite ratio of conversion between the various special forces of
nature (motion, heat, light, electricity, etc.), and it was
now carried through by experimental investigations. It now stands as
a law for all material manifestations of force. --- Nowhere is it
so clear as here that the science of the nineteenth century, in all its
greatness, consists after all in carrying through ideas that stem from
the remarkable time when modern science was founded.

It has been said of the principle of the conservation of energy that it
is only a generalization from a limited experience. It is also
true that a confirmation by experience of a principle such as this can in a certain
sense never be concluded. One will never be able to isolate and
close off a part of the world in such a way that every supply or removal
of energy would be absolutely excluded. And there always remains the
possibility that the principle holds only for the part or parts
of the world that we have so far been able to investigate. And since ``the world'' can never
come to lie before us as a whole concluded once and for all,
it is an infinite task that is set for research here. But the
% --- p. 29 ---
% running head: ROBERT MAYER OG SADI CARNOT (Robert Mayer and Sadi Carnot)
princi\opage{29}ple of the conservation of energy here only shares the fate of all the
propositions to which research into actual nature can lead. The
concept ``world'' in the sense of absolute totality has the scientific
significance that it shows us the possibility of ever new tasks
as the horizon widens. And it marks a significant
advance that the principle put forward in the seventeenth century has now
proved accessible to progressive confirmation by experience. Like
every great hypothesis, the principle supplies an important maxim for
research, namely, always to look for new forms of energy as
equivalents for those that disappeared, and to look for the preceding
forms of energy from which the new ones have arisen. Maxwell has
particularly emphasized the epistemological significance of the principle
of the conservation of energy, that it leads us, at every new
natural phenomenon, to put definite questions. This is indeed the significance
of all living principles in science, and this significance
cannot be overstated.

The principle says that in every material change there is only a conversion
of energy, without any energy being lost; but it
says nothing about the direction in which the conversions take place, and
in particular nothing about whether a conversion in the opposite direction can always as
easily take place. A French investigator had, as early as eighteen years before
Mayer's treatise, put forward a principle that throws light on this matter.
Sadi Carnot (\textit{Réflexions sur la puissance motrice du feu}. 1824)
showed that, for heat to be a motive force, there
must be a difference of temperature. Motion is produced only
when such a difference arises between the states of the bodies
between which the action is to take place. But since temperature
always tends toward equilibrium, heat passing over from warmer
to colder bodies and thereby being distributed, the possibility
of a new conversion of heat into motion thereby falls away within
the same circle of bodies, considered as a closed whole.

Carnot's principle, too, holds only for a limited world. And
if by ``the world'' one understands a limited whole, there prevails in
it a process that leads in the direction of a complete equilibrium, a
state in which no new release of energy is possible.

Carnot's principle, which, like Mayer's, had to wait long for
full recognition, does not conflict with the principle of the conservation of
energy. For in the state of equilibrium that there is a constant tendency to bring
about, no energy has been lost; only, all energy has passed over from
being actual to being potential.

% --- p. 30 ---
% running head: HORIZONTEN UDVIDES NEDAD OG UDAD (The horizon widens downward and outward)
\opage{30}An optimistic philosophy of nature will hold to Mayer's principle, a
pessimistic one to Carnot's. But both principles hold only for a
limited horizon, and since science always seeks to widen
its horizon, the discussion of which standpoint is the decisive
one will always have to be taken up anew. ---

A still more intimate connection between the various forces of
nature than that which follows from the principle of the conservation of energy
was shown when it was established that forces of nature which had so far stood as
different in kind follow the same laws and must therefore be assumed to be
of the same nature. Thus Melloni and other investigators showed that
light and heat were one. H. C. Ørsted had already divined a connection
between light and electricity, conceiving the lightning flash as
a series of electric sparks. And after Maxwell had by a mathematical route
put forward the electromagnetic theory of light, according to which light
can be traced back to electricity, the laws of these two activities
being the same, Hertz showed experimentally that light and electricity,
as regards wave motion, velocity, refraction, etc.,
present the same relations. Optics could now be regarded as an appendix
to the theory of electricity.

On the whole, the confirmation by experience that the principle
of the conservation of matter and force gained in the nineteenth century
is far from having brought the scientific conception of the world to
a conclusion. It was in itself a magnificent picture of the world that was unrolled.
The world came to stand as a whole within which constant
conversions of matter and force took place, without any matter or force at all
arising or perishing. As we have already seen, this picture of
the world cannot be concluded, since the widening of the horizon can always
bring new possibilities. But on another side too, in the very foundation,
research goes beyond the picture of the world just indicated. The
development of the theory of electricity has led to this: that far down beneath the
matters and forces of earlier science a new and rich world comes
into view, in which perhaps the real foundation of the world of matter
and force in which we move is to be sought. The theory of electricity
has led to the discovery of minute particles, so-called electrons, which are
smaller than Dalton's atoms and of which these must be thought to be
composed. The change in the scientific conception that was thereby
heralded at the end of the century is described in the following
way by a Danish physicist: ``In the last half of the century
the theory of electricity has pressed forward so strongly at the expense of the other
forces of nature that it seems about to swallow them
% --- p. 31 ---
% running head: EN NY FYSIK (A new physics)
\opage{31}all. One is now inclined in advance to refer not only
optics but also the other forces of nature to it; electricity
alone is so manifold in its forms of manifestation that it can be thought
capable of comprising them all. Even the very concept of mass can be
dispensed with, if the atoms are thought of as charged with electricity, since
in that case every acceleration must produce induction in the surroundings,
and thereby in turn a resistance to the motion. The dissolving power of
the theory of electricity has extended to chemistry as well, to which
physical chemistry bears witness. Even the old concept of the atom
is in danger of collapsing. Thus J. J. Thomson and
Zeemann have believed they could demonstrate masses a thousand times lighter
than the chemists' atoms'' (C. Christiansen: ``Før og nu'' [Then and now], first
volume of \textit{Tidsskrift for Fysik}. 1908). % Zeemann as printed (for Zeeman)

There thus opens up, at the end of the century, the prospect of
an entirely new physics, which goes further back than the one founded by Galilei and Newton
and will derive the maxims of the old physics from still
deeper-lying principles. From the world of the great masses one has
begun to descend to the world of the small elements. ``The time,'' says
an American physicist, ``seems not to be distant when there will be
room for a new `Principia' and for a work that will accomplish for the
molecular systems what the `Mécanique Céliste' accomplished for the solar
system.'' (Woodward: The Unity of Physical Science. Congress of
Arts and Sciences St. Louis 1904). % Woodward translated back from Høffding's Danish; English original not checked.  Céliste as printed.

If it is asked what it is, then, that moves according to the laws of electricity
and under special conditions appears as light, heat,
etc., the natural scientist answers with the beautiful name ether. But
the ether is known only through its special forms of manifestation, among
which electricity is regarded as the fundamental one. The famous
French physicist Langevin declares that it is impossible to give a picture,
a representation, of the ether, since its properties are entirely different
from those of matter and probably far simpler than these.
And he thinks that the difficulty we face here perhaps
hangs together with the present constitution of human consciousness:
for centuries we have been used to thinking by means
of the concept of matter, and are therefore unable to form a
concrete representation that does not have a material content. The ether
can be thought only by means of the concepts of the electromagnetic
properties we ascribe to it (Langevin: Congress of Arts and
Sciences. 1904).

Plato already saw that a thoroughgoing atomic theory cannot be built
% --- p. 32 ---
% running head: MEKANISME OG VITALISME (Mechanism and vitalism)
\opage{32}up by means of the ordinary concept of matter, cannot be grounded in
sensuous representations, but must make use of analogy. The result
he reached by a critical thinking-through of Democritus's atomic
theory, modern physics has now reached by its own route. ---

For no science have greater possibilities and tasks shown themselves at the beginning
of the twentieth century than for physics. But
what has already been won --- the progressive carrying-through and
confirmation of the great ideas that three hundred years ago led
to the rise of a ``new science'' --- is a great achievement, which will be of
incalculable significance for the human conception of existence,
once its scope really dawns on consciousness.
Both in the great and in the small, in the world of the heavenly bodies and in that of the
imperceptible elements, a greater connection has shown itself
than one could have divined.

\sectionhead{b.}{The Theory of Evolution.}

The science of organic life could not remain untouched
by the great movement stirring in the physical and chemical
domain. Ever since the first rise of modern natural science, one had
sought to bring the doctrine of life under the general laws of
nature as physics demonstrates them. The principle that all change
is motion, one sought to maintain not only in the domain of inorganic
but also in that of organic nature. Descartes
insisted emphatically that organisms were only mechanisms,
set going according to the general laws of motion, and he
held that Harvey's discovery of the circulation of the blood confirmed this
view, which became the prevailing one in the seventeenth century. As a
reaction against this mechanical view there appeared the so-called
vitalism, which maintained that organic phenomena could find
their explanation only if one assumed a peculiar life force, independent of all physical and
chemical forces. Vitalism reached its most pronounced form
in the medical school of Montpellier in the eighteenth century.
The ``life force'' directs the physical and chemical processes in the organism in
definite directions, brings about a connection among everything that
goes on in the living being, so that it can subsist as a whole,
provides for the formation of the substances that life needs, and for sufficient
strength to overcome the resistance that life meets from without.

The nineteenth century appears here, as on other points,
as a carrying-through and specialization of the fundamental ideas of the seventeenth century.
Mechanism and vitalism at bottom designate only two
[\dots]

\end{document}
